\documentclass[12pt,a4paper]{article}

% =========================
% Pacotes essenciais
% =========================
\usepackage[a4paper,margin=2.5cm]{geometry}
\usepackage[T1]{fontenc}
\usepackage[utf8]{inputenc}
\usepackage[brazil]{babel}
\usepackage{lmodern}
\usepackage{setspace}
\usepackage{microtype}
\usepackage{hyperref}
\usepackage{enumitem}
\usepackage{fancyhdr}

\hypersetup{
  colorlinks=true,
  linkcolor=black,
  urlcolor=blue,
  citecolor=black
}

\setstretch{1.15}
\setlength{\parindent}{0pt}
\setlength{\parskip}{6pt}

\pagestyle{fancy}
\fancyhf{}
\lhead{Relatório do Piloto — CV721A}
\rhead{\thepage}
\cfoot{}

% =========================
% Metadados
% =========================
\newcommand{\Disciplina}{Fundações (CV721A)}
\newcommand{\Curso}{Engenharia Civil — UNICAMP}
\newcommand{\Periodo}{2025}
\newcommand{\VersaoRelatorio}{v1.0}
\newcommand{\DataRelatorio}{\today}

\title{\textbf{Relatório do Projeto Piloto}\\
Sistema de Apoio ao Trabalho Docente\\
\Disciplina}
\author{}
\date{\DataRelatorio\ (\VersaoRelatorio)}

\begin{document}
\maketitle

% =========================================================
\section{Introdução}

O presente relatório descreve o desenvolvimento e a aplicação de um \textbf{Sistema de Apoio ao Trabalho Docente}, testado como projeto piloto na disciplina \Disciplina{}, do curso de \Curso{}, no período letivo de \Periodo{}.

O objetivo principal do projeto é reduzir a carga operacional associada ao controle de presença e à consolidação de notas, permitindo que o professor concentre esforços na condução pedagógica da disciplina.

% =========================================================
\section{Contextualização do Piloto}

Em disciplinas com grande número de alunos, tarefas como registro manual de presença, consolidação de resultados de atividades interativas e lançamento de notas tornam-se processos repetitivos, suscetíveis a erros e consumidores de tempo.

Ferramentas institucionais existentes, embora essenciais, nem sempre são otimizadas para o fluxo diário da aula presencial. Nesse contexto, propôs-se o desenvolvimento de uma ferramenta complementar, centrada no professor, que se integre de forma leve às práticas já consolidadas.

% =========================================================
\section{Descrição do Sistema Desenvolvido}

O sistema desenvolvido consiste em uma aplicação web acessada via navegador, organizada em um painel único para o professor, com os seguintes módulos principais:

\begin{itemize}[noitemsep]
  \item Aula de Hoje;
  \item Controle de Frequência;
  \item Notas e Avaliações;
  \item Importação de resultados do Kahoot;
  \item Exportação de dados para o sistema institucional.
\end{itemize}

O controle de presença é realizado por meio de QR Code com validade temporária, permitindo o registro automático da presença dos alunos a partir de seus dispositivos móveis.

As notas podem ser inseridas manualmente ou importadas a partir de arquivos CSV gerados por plataformas de atividades interativas, como o Kahoot.

% =========================================================
\section{Metodologia do Piloto}

O projeto foi conduzido como um estudo piloto aplicado a uma turma real da disciplina \Disciplina{}.

A metodologia adotada incluiu:
\begin{itemize}[noitemsep]
  \item desenvolvimento incremental do sistema;
  \item validação funcional em ambiente real de sala de aula;
  \item observação do uso pelo professor responsável;
  \item coleta de percepções qualitativas sobre usabilidade e utilidade.
\end{itemize}

Não houve substituição de sistemas institucionais existentes, sendo o sistema utilizado de forma complementar.

% =========================================================
\section{Resultados Preliminares}

Durante o período inicial de testes, observou-se:
\begin{itemize}[noitemsep]
  \item redução do tempo necessário para registro de presença;
  \item eliminação de registros manuais em papel;
  \item maior organização dos dados de frequência;
  \item facilidade na incorporação de atividades interativas como instrumento avaliativo.
\end{itemize}

O uso do QR Code mostrou-se adequado ao contexto de sala de aula, com boa aceitação por parte dos alunos.

% =========================================================
\section{Discussão}

Os resultados preliminares indicam que o sistema atende ao objetivo de apoiar o trabalho docente sem introduzir complexidade excessiva.

A adoção de uma abordagem incremental permitiu ajustes rápidos conforme as necessidades observadas em sala, reforçando a importância de soluções flexíveis no contexto do ensino superior.

A integração leve com ferramentas já utilizadas pelo professor mostrou-se um fator determinante para a aceitação do sistema.

% =========================================================
\section{Considerações Finais}

O projeto piloto demonstrou a viabilidade de um sistema dedicado ao apoio ao trabalho docente, focado na automação de tarefas operacionais e na integração com práticas pedagógicas existentes.

O sistema apresenta potencial para expansão e adaptação a outras disciplinas, respeitando as especificidades de cada contexto.

% =========================================================
\section{Trabalhos Futuros}

Como trabalhos futuros, destacam-se:
\begin{itemize}[noitemsep]
  \item ampliação dos módulos de análise e visualização de dados;
  \item integração mais direta com sistemas institucionais;
  \item avaliação quantitativa do impacto pedagógico;
  \item extensão do piloto para outras disciplinas.
\end{itemize}

\end{document}
